\chapter{Allophony}

Certain other changes occur purely for the purposes of euphony or minimising
articulatory effort, and therefore are not phonemic in nature.
The changes described in this section are thus not represented in writing.

\section{Monophthongisation}

The diphthongs \phonem{ɯj, ij, uw, yw} are subject to a process of
monophthongisation, wherein the glide is dropped and the vowel itself is
lengthened.

\ex
  \vtop{\halign{
    \phonem{#}\hfil&\:\alternate~\phonet{#}\hfil&\:\transl{#}\cr
    kij-&kiː&put on, don\cr
    suw&suː&water\cr
  }}
\xe

\section{Nasal assimilations}

Before \phonem{G}, \phonem{n} may be realised as \phonet{ŋ}~(\ref{%
ex:regressive-velarisation}).
After nasal consonants, \phonem{l} may be realised as  \phonet{n}~(\ref{%
ex:progressive-nasalisation}).

\pex
  \a\label{ex:regressive-velarisation}
    \vtop{\halign{
      \phonem{#}\hfil&\:\alternate~\phonet{#}\hfil&\:\transl{#}\cr
      d͡ʒɛng&d͡ʒɛŋg&war\cr
      sɛngɛɾ&sɛŋgɛɾ&trench, dugout\cr
    }}
  \a\label{ex:progressive-nasalisation}
    \vtop{\halign{
      \phonem{#}\hfil&\:\alternate~\phonet{#}\hfil&\:\transl{#}\cr
      qumlɑɾ&qumnɑɾ&sands\cr
      ɔɾmɑnlɑɾ&ɔɾmɑnnɑɾ&forests\cr
      ɑŋlɑmɑq&ɑŋnɑmɑq&to understand\cr
    }}
\xe
